%------------------------------------------------------------
% CV – Daniele Moretto
% Style inspired by Nicolò Crescenzio's CV
% Requires: geometry, titlesec, enumitem, hyperref, fontawesome5,
%           xcolor, parskip, array, tabularx, microtype
%------------------------------------------------------------
\documentclass[10pt, a4paper]{article}
 
% ── Packages ────────────────────────────────────────────────
\usepackage[
  top=1.8cm, bottom=1.8cm,
  left=2.0cm, right=2.0cm
]{geometry}
\usepackage{fontawesome5}          % icons
\usepackage{xcolor}
\usepackage{titlesec}
\usepackage{enumitem}
\usepackage{hyperref}
\usepackage{array}
\usepackage{tabularx}
\usepackage{microtype}
\usepackage{parskip}
\usepackage{lmodern}
\usepackage[T1]{fontenc}
\usepackage[utf8]{inputenc}
 
% ── Colours ─────────────────────────────────────────────────
\definecolor{accent}{HTML}{2E4057}   % dark slate blue
\definecolor{lightgray}{HTML}{666666}
 
% ── Hyperref setup ──────────────────────────────────────────
\hypersetup{
  colorlinks = true,
  urlcolor   = accent,
  linkcolor  = accent
}
 
% ── Section headings ────────────────────────────────────────
\titleformat{\section}
  {\large\bfseries\color{accent}}
  {}{0em}{}
  [\vspace{-0.4em}\textcolor{accent}{\rule{\linewidth}{0.6pt}}\vspace{0.2em}]
 
\titlespacing{\section}{0pt}{1.2em}{0.4em}
 
% ── No page numbers ─────────────────────────────────────────
\pagestyle{empty}
 
% ── Timeline entry helper ───────────────────────────────────
% \entry{date range}{Title}{Institution}{Description bullets}
\newcommand{\entry}[4]{%
  \begin{tabular}{@{}p{3.2cm} p{12.0cm}@{}}
    \raggedright\small\textcolor{lightgray}{#1} &
    \begin{minipage}[t]{12.0cm}
      \textbf{#2}\\
      \textit{#3}
      #4
    \end{minipage}
  \end{tabular}%
  \vspace{0.55em}
}
 
% ── Compact bullet list inside entries ──────────────────────
\newenvironment{entrylist}{%
  \begin{itemize}[leftmargin=1em, itemsep=0pt, topsep=2pt, parsep=0pt]
}{%
  \end{itemize}
}
 
%============================================================
\begin{document}
 
%------------------------------------------------------------
% HEADER
%------------------------------------------------------------
\begin{center}
  {\LARGE\bfseries Daniele Moretto}\\[0.5em]
  \textcolor{lightgray}{
    \faMapMarker*\enspace Via Marzolo, 9 - 35131 Padova (PD), Italy
    \quad
    \faPhone*\enspace +39 328 3763380
  }\\[0.3em]
  \href{mailto:daniele.moretto.3@phd.unipd.it}{%
    \faEnvelope\enspace daniele.moretto.3@phd.unipd.it}
\end{center}
 
\vspace{0.6em}
 
%------------------------------------------------------------
% ABOUT
%------------------------------------------------------------
I am a PhD student in Sciences of Civil, Environmental and Architectural Engineering
at the University of Padova. My research is focused on numerical methods for coupled
multi-physics problems in porous media, reservoir simulation, and the development of
general-purpose computational frameworks for subsurface applications.
 
%------------------------------------------------------------
% EDUCATION
%------------------------------------------------------------
\section{Education}
 
\entry{Oct 2023 – Present}
  {PhD in Sciences of Civil, Environmental and Architectural Engineering}
  {Università degli Studi di Padova}
  {%
    \begin{entrylist}
      \item Advisors: Prof.\ Massimiliano Ferronato and Prof.\ Andrea Franceschini.
    \end{entrylist}
  }
 
\entry{2021 – 2023}
  {Master's Degree in Civil Engineering}
  {Università degli Studi di Padova}
  {%
    \begin{entrylist}
      \item Thesis: \textit{GReS – General Reservoir Simulator: development of a
        general-purpose simulator for coupled multi-physics reservoir problems}.
      \item Advisors: Prof.\ Massimiliano Ferronato, Prof.\ Andrea Franceschini.
      \item Grade: 110/110 \textit{cum laude}.
    \end{entrylist}
  }
 
\entry{2018 – 2021}
  {Bachelor's Degree in Civil Engineering}
  {Università degli Studi di Padova}
  {%
    \begin{entrylist}
      \item Thesis: \textit{Finite element method for the solution of the
        mechanical problem in porous media}.
      \item Advisor: Prof.\ Massimiliano Ferronato.
      \item Grade: 110/110 \textit{cum laude}.
    \end{entrylist}
  }
 
%------------------------------------------------------------
% RESEARCH & WORK EXPERIENCE
%------------------------------------------------------------
\section{Research \& Teaching Experience}
 
\entry{Oct 2023 – Present}
  {Doctor of Philosophy (PhD)}
  {Università degli Studi di Padova – DICEA}
  {%
    \begin{entrylist}
      \item Research on numerical algorithms for coupled multiphysics and multidomain simulations in fractured porous media
    \end{entrylist}
  }
 
\entry{2024}
  {Integrated Didactic – Numerical Analysis (laboratory)}
  {Università degli Studi di Padova}
  {%
    \begin{entrylist}
      \item Teaching assistant for undergraduate numerical analysis laboratory.
    \end{entrylist}
  }
 
\entry{2023}
  {Integrated Didactic – Calculus (exercises)}
  {Università degli Studi di Padova}
  {%
    \begin{entrylist}
      \item Teaching assistant for undergraduate calculus course.
    \end{entrylist}
  }
 
%------------------------------------------------------------
% JOURNAL PAPERS
%------------------------------------------------------------

 \input{publications}
 

 
%------------------------------------------------------------
% COMPUTER SKILLS
%------------------------------------------------------------
\section{Computer Skills}
 
\renewcommand{\arraystretch}{1.3}
\begin{tabular}{@{} p{6cm} p{12.4cm} @{}}
  \textbf{Coding}     & C++, Python, MATLAB, Bash, \LaTeX \\
  \textbf{Software and other tools}   & ParaView, Abaqus, Microsoft Office \\
\end{tabular}
 
%------------------------------------------------------------
% LANGUAGES
%------------------------------------------------------------
\section{Languages}
 
\renewcommand{\arraystretch}{1.3}
\begin{tabular}{@{} p{2.8cm} p{12.4cm} @{}}
  \textbf{Italian} & Mother Tongue \\
  \textbf{English} & TOEFL C1 level \\
\end{tabular}
 
\end{document}